%-------------------------
% Cover Letter - NetApp Emerging Talent Entry Level SE (Systems)
% Author: Ajay Venkatesh
%-------------------------

\documentclass[letterpaper,11pt]{article}

\usepackage[top=0.8in, bottom=0.8in, left=0.9in, right=0.9in]{geometry}
\usepackage[hidelinks]{hyperref}
\usepackage{lmodern}
\usepackage{parskip}
\usepackage{microtype}

\setlength{\parskip}{6pt}
\setlength{\parindent}{0pt}

\begin{document}

\begin{flushleft}
\textbf{\large Ajay Venkatesh} \\
Brooklyn, NY \\
+1 516 585 9013 \\
\href{mailto:av3855@nyu.edu}{av3855@nyu.edu}
\end{flushleft}

\vspace{10pt}

May 15, 2026

\vspace{6pt}

Hiring Manager \\
NetApp

\vspace{10pt}

Dear Hiring Manager,

My name is Ajay Venkatesh --- finishing my M.S. in Computer Engineering at NYU, with several years of production software experience: a few years at PwC in Bengaluru, a summer SDE internship at Amazon, and ongoing on-campus work at NYU's Novel AI Technologies. My graduate work has been weighted toward the lower end of the stack --- \textbf{Computer Architecture, System Design, Data Structures,} and \textbf{Big Data} alongside ML and Deep Learning --- which is the same layer-cake an ONTAP role lives in.

The Emerging Talent posting maps onto where I've been pushing my own learning. \textbf{C++} is where I cut my teeth and where I keep going back when an abstraction stops being honest about what's happening underneath. I've practiced multithreaded programming through coursework and in projects --- a Python desktop application I shipped handled concurrent transaction processing and real-time cost splitting using threading and synchronized locks for thread-safe persistence. The shape of that work --- concurrency correctness, careful resource management, deterministic behavior under contention --- is the same shape systems software demands.

On engineering discipline: at Amazon I engineered backend services in \textbf{Java} (Coral, Smithy, Dagger, CBOR) with low-level design patterns and unit/integration testing via \textbf{Mockito}, shipped under code-review discipline. I treat tests as part of the build, not an afterthought, and I'm comfortable navigating large codebases, debugging through logs, metrics, and traces, and the actual unglamorous work of keeping software reliable. At PwC I led code reviews and mentored engineers on software design principles, contributing to a measurable drop in production defects.

On collaboration and the product development cycle: at PwC I spent three years partnering with architects, product managers, and business stakeholders on enterprise systems --- definition, design, implementation, testing, deployment, and early customer support across the full lifecycle. That cross-functional posture, plus the customer-deployment understanding the role calls out, is something I've already been practicing.

What pulls me toward NetApp is the seriousness of the product space. Storage operating systems sit at a layer where correctness and performance both matter in real units, where bugs have real consequences for customers' data, and where engineering judgment has to hold up under decades of compatibility constraints. That's the kind of environment I want to grow in, and ONTAP is one of the most respected systems in the field to grow inside of.

I'd welcome the chance to discuss further. Thanks for considering my application.

\vspace{12pt}

Sincerely, \\
Ajay Venkatesh

\end{document}
